\section*{第 \ref{ch:g5:raw-materials-and-recycling} 章——从原料到回收利用}

\begin{solution}{exo:g5:raw-materials-and-recycling:1}
\omterm{def:g5:raw-materials-and-recycling:natural-material}{天然材料}：木材、羊毛、石头。\omterm{def:g5:raw-materials-and-recycling:natural-material}{人造材料}：玻璃、钢、塑料。
\end{solution}

\begin{solution}{exo:g5:raw-materials-and-recycling:2}
纸：木材（树）。玻璃：沙子（和苏打、石灰石一起）。
棉花：棉花植株。塑料：原油。
\end{solution}

\begin{solution}{exo:g5:raw-materials-and-recycling:3}
从中提取金属的岩石。铝取自铝土矿。
\end{solution}

\begin{solution}{exo:g5:raw-materials-and-recycling:4}
玻璃罐放入玻璃回收箱，报纸和纸板放入纸类回收箱，
金属罐放入金属回收箱，塑料瓶放入塑料回收箱。
\end{solution}

\begin{solution}{exo:g5:raw-materials-and-recycling:5}
新原料：用回收材料代替取自自然界的沙子、\omterm{def:g5:raw-materials-and-recycling:raw-material}{矿石}或
原油。
\end{solution}

\begin{solution}{exo:g5:raw-materials-and-recycling:6}
可再生：木材（若重新种植）、羊毛、棉花。不可再生：原油、
铁\omterm{def:g5:raw-materials-and-recycling:raw-material}{矿石}。
\end{solution}

\begin{solution}{exo:g5:raw-materials-and-recycling:7}
用沙子制造玻璃是\omterm{def:g5:irreversible-changes:chemical-change}{化学变化}：生成了新物质，
且无法简单地恢复原样。
\end{solution}

\begin{solution}{exo:g5:raw-materials-and-recycling:8}
例如：减少——购买不带塑料包装的水果；重复使用——
重新灌装玻璃罐；\omterm{def:g5:raw-materials-and-recycling:recycling}{回收利用}——把饮料罐放入金属回收箱。
\end{solution}

\begin{solution}{exo:g5:raw-materials-and-recycling:9}
每种材料都有各自的回收方式：混入其他材料的碎片
会破坏熔化的玻璃。
\end{solution}

\begin{solution}{exo:g5:raw-materials-and-recycling:10}
20 单位的二十分之一：$20 \div 20 = 1$ 单位能量（大约）。
\end{solution}

\begin{solution}{exo:g5:raw-materials-and-recycling:11}
\omterm{def:g5:raw-materials-and-recycling:recycling}{回收利用}仍需要能量、机器和运输；
不被生产出来的物品，能节省原料和上述全部能量。
\end{solution}
